\hypertarget{the-six-birds-vocabulary}{%
\section{The Six Birds vocabulary}\label{the-six-birds-vocabulary}}

This section introduces only the parts of the Six Birds closure calculus
that the paper uses. Fuller accounts are in \citep{TsiokosSixBirds2026}
and \citep{TsiokosFoundations2026}.

\hypertarget{theories-as-closure-packages}{%
\subsection{Theories as closure
packages}\label{theories-as-closure-packages}}

The calculus describes theories by how they compress and complete a
system. A \textbf{closure package} consists of a \textbf{lens} \(f\), a
deterministic coarse-graining that keeps some features of a substrate
and discards the rest; a \textbf{completion} \(E\), which fills in what
the lens discarded whenever the description is used and is idempotent on
its declared domain (\(E^2=E\)); and \textbf{audits}, functionals that
test whether the package stays consistent in use. The standard
cosmological model can be read in these terms: the FLRW lens keeps a few
homogeneous degrees of freedom, the homogeneous completion fills in the
rest, and consistency tests between probes act as audits. We use this
reading as a guide for building models, not as a theorem about
cosmology.

The program organizes the moves a package can make into six primitive
roles, \(P_1\) to \(P_6\). The Foundations series argues that these six
are canonical under its composability and bounded-access assumptions
\citep{TsiokosFoundations2026}; that argument is not repeated or relied
on here. The roles that appear in this paper are:

\begin{longtable}[]{@{}lll@{}}
\toprule
\begin{minipage}[b]{0.05\columnwidth}\raggedright
role\strut
\end{minipage} & \begin{minipage}[b]{0.26\columnwidth}\raggedright
meaning\strut
\end{minipage} & \begin{minipage}[b]{0.60\columnwidth}\raggedright
where it appears\strut
\end{minipage}\tabularnewline
\midrule
\endhead
\begin{minipage}[t]{0.05\columnwidth}\raggedright
\(P_1\)\strut
\end{minipage} & \begin{minipage}[t]{0.26\columnwidth}\raggedright
operator rewrite\strut
\end{minipage} & \begin{minipage}[t]{0.60\columnwidth}\raggedright
a correction term added to the macro law, such as the rewrite family
\(\Delta_H(z)\) (Section 8)\strut
\end{minipage}\tabularnewline
\begin{minipage}[t]{0.05\columnwidth}\raggedright
\(P_2\)\strut
\end{minipage} & \begin{minipage}[t]{0.26\columnwidth}\raggedright
constraints and gating\strut
\end{minipage} & \begin{minipage}[t]{0.60\columnwidth}\raggedright
admissibility rules: budgets, collapse thresholds, convergence tests
(Sections 4--6)\strut
\end{minipage}\tabularnewline
\begin{minipage}[t]{0.05\columnwidth}\raggedright
\(P_3\)\strut
\end{minipage} & \begin{minipage}[t]{0.26\columnwidth}\raggedright
protocols and route structure\strut
\end{minipage} & \begin{minipage}[t]{0.60\columnwidth}\raggedright
the order in which completions are applied, and probe routes in the data
(Sections 6, 8)\strut
\end{minipage}\tabularnewline
\begin{minipage}[t]{0.05\columnwidth}\raggedright
\(P_4\)\strut
\end{minipage} & \begin{minipage}[t]{0.26\columnwidth}\raggedright
staging and refinement\strut
\end{minipage} & \begin{minipage}[t]{0.60\columnwidth}\raggedright
redshift bins, probe splits, finite iteration budgets (Sections 6,
8)\strut
\end{minipage}\tabularnewline
\begin{minipage}[t]{0.05\columnwidth}\raggedright
\(P_5\)\strut
\end{minipage} & \begin{minipage}[t]{0.26\columnwidth}\raggedright
packaging\strut
\end{minipage} & \begin{minipage}[t]{0.60\columnwidth}\raggedright
the quotient into macro-objects: patches, records, the homogeneous
background\strut
\end{minipage}\tabularnewline
\begin{minipage}[t]{0.05\columnwidth}\raggedright
\(P_6\)\strut
\end{minipage} & \begin{minipage}[t]{0.26\columnwidth}\raggedright
accounting\strut
\end{minipage} & \begin{minipage}[t]{0.60\columnwidth}\raggedright
budgets and ledgers, whose multipliers are the shadow prices of Section
4\strut
\end{minipage}\tabularnewline
\bottomrule
\end{longtable}
\addtocounter{table}{-1}

\hypertarget{readouts-and-determination}{%
\subsection{Readouts and
determination}\label{readouts-and-determination}}

Many results in this paper answer the question \emph{does this record
determine that quantity?} Let \(r:Z\to Y\) be a readout on a nonempty
set \(Z\) of admissible configurations and \(q:Z\to Q\) a quantity of
interest. We say that \(r\) \textbf{determines} \(q\) if
\(q=\phi\circ r\) for some function \(\phi:Y\to Q\), that is, if
configurations with equal readouts always have equal values of \(q\). If
two admissible configurations have the same readout but different values
of \(q\), no decoder, however complicated, can recover \(q\) from \(r\).
This is the factorization test of the calculus. It is a statement about
the readout, not about all possible observations: a richer readout may
determine \(q\).

A second test asks whether two completions commute, \(E_1E_2=E_2E_1\). A
nonzero commutator means that the order in which two fill-in rules are
applied matters. Section 6.5 applies both tests to a pair of coarse
descriptions.

\hypertarget{budgets-and-shadow-prices}{%
\subsection{Budgets and shadow
prices}\label{budgets-and-shadow-prices}}

The calculus records a conserved or budgeted quantity at one level as a
\textbf{price} at the level above. If a description minimizes a cost
\(J\) subject to a budget \(e(x)\le B\), then under convexity and
regularity conditions (Section 4.1) the multiplier of the budget is the
rate at which the optimal cost \(V(B)\) falls as the budget is relaxed,
\(p=-V'(B)\). Inverse temperature as the multiplier of a mean-energy
constraint in entropy maximization, and prices in resource allocation,
are familiar examples. What such a construction fixes is the
\emph{existence} and the \emph{form} of the multiplier: one number per
binding global constraint, zero when the constraint is slack. Its
numerical value depends on the budget level and on the objective.
Section 4 states the precise theorem used here.

In the Six Birds corpus this reading was also applied to holographic
min-cut duality \citep{RyuTakayanagi2006, FreedmanHeadrick2017}. The
precise finite relation, which we use only as a comparison, is stated in
Section 4.4.

\hypertarget{evidence-types-and-provenance}{%
\subsection{Evidence types and
provenance}\label{evidence-types-and-provenance}}

We report every result with its evidence type:

\begin{itemize}
\tightlist
\item
  \textbf{Theorem {[}Lean{]}}: a mathematical statement about a declared
  model whose proof is checked by the Lean 4 kernel with mathlib
  \citep{deMoura2021, mathlib2020}, using only the standard axioms
  \texttt{propext}, \texttt{Classical.choice} and \texttt{Quot.sound}.
  Statements that rely on an argument not yet mechanized are labelled as
  such.
\item
  \textbf{Exact computation}: a computation in rational or interval
  arithmetic outside Lean, with its inputs recorded.
\item
  \textbf{Floating computation}: an ordinary numerical computation. Its
  reported precision is not a rigorous error bound.
\item
  \textbf{Audit}: a pre-registered statistical test on public data, with
  its calibration and any amendments disclosed.
\end{itemize}

Results computed in this paper's repository are called \textbf{local}.
Results taken from other Six Birds papers are called \textbf{imported};
they motivate the models and are never used as local proofs. One
imported computation, the background sweep behind Section 6, was copied
into the repository with a hash manifest, re-run, and re-derived by
independent code; its rows are therefore local results, and Section 6
gives their mathematical analysis. Several cross-paper connections drawn
elsewhere in the Six Birds corpus are not used, because their current
source does not establish them as theorems: a blindness of the calculus
to the Standard Model's generation count, an exclusion of proton decay
and monopoles by record stability, and a null prediction for tabletop
tests of gravitationally induced entanglement, which that source derives
only for an assumed perfect-record channel.
